\lstset{literate=
{π}{{$\pi$}}1
{σ}{{$\sigma$}}1
{γ}{{$\gamma$}}1
{ρ}{{$\rho$}}1
{∩}{{$\cap$}}1
{←}{{$\leftarrow$}}1
{→}{{$\to$}}1
{⟗}{{$\Join_{\mathrm{full}}$}}1
{⨯}{{$\times$}}1
{⨝}{{$\Join$}}1
{∪}{{$\cup$}}1
{⋈}{{$\bowtie$}}1
{∧}{{$\land$}}1
{∨}{{$\lor$}}1
{≥}{{$\geq$}}1
{≤}{{$\leq$}}1
{⟕}{{$\leftouterjoin$}}1
{⟖}{{$\rightouterjoin$}}1
}

\seccion{1. Cardinalidad de las consultas}

Las relaciones iniciales son:
\[
R(A,B)=\{(1,x),(2,y),(2,z),(3,x),(9,a)\},\qquad
S(B,C,D)=\{(x,0,3),(y,2,1),(y,3,3),(w,3,0),(y,4,2)\}.
\]

Se usa semántica de conjuntos. Las proyecciones eliminan las tuplas duplicadas. En C10, el Join externo completo conserva las tuplas que no encuentran coincidencia y completa con \nulo{} los atributos faltantes.

\subsection*{C08. Intersección de proyecciones}

Primero se proyecta el atributo $A$ de $R$:
\[
\pi_A(R)=\{1,2,3,9\}.
\]

Aunque $R$ contiene dos tuplas con $A=2$, la proyección conserva ese valor una sola vez. Después se proyecta $D$ de $S$:
\[
\pi_D(S)=\{3,1,0,2\}.
\]

El valor $3$ aparece dos veces en $S$, pero se conserva una sola vez. Se renombra $D$ a $A$ y se calcula la intersección:
\[
\rho_{A\leftarrow D}(\pi_D(S))=\{3,1,0,2\},\qquad
\{1,2,3,9\}\cap\{3,1,0,2\}=\{1,2,3\}.
\]

\begin{center}
\begin{tabular}{c}
\toprule
\rowcolor{Fondo}
$A$\\
\midrule
1\\
2\\
3\\
\bottomrule
\end{tabular}
\end{center}

\textbf{Cardinalidad: 3.} Los valores comunes son 1, 2 y 3. El valor 9 no pertenece a la segunda relación, mientras que el 0 no pertenece a la primera.

\subsection*{C10. Join externo completo y agrupación}

EL join externo completo combina las tuplas que coinciden en $B$ y conserva las que no tienen pareja. En la expresión ejecutada en RelaX, se renombra el atributo $B$ de $S$ como $BS$ para evitar nombres de columna duplicados. Después, \texttt{coalesce} permite recuperar un único atributo $B$, tomando el valor no nulo de cualquiera de las relaciones.

\[
S_2\leftarrow\rho_{S_2}(\rho_{BS\leftarrow B}(S)).
\]

\begin{lstlisting}
J = π A, coalesce(B,BS) → B, C, D (R ⟗ B = BS S2)
\end{lstlisting}


Las coincidencias son: $(1,x)$ con la tupla de $S$ cuyo $B=x$; $(2,y)$ con las tres tuplas de $S$ cuyo $B=y$; y $(3,x)$ con la tupla de $S$ cuyo $B=x$. También se conservan las tuplas de $R$ sin coincidencia, $(2,z)$ y $(9,a)$, así como la tupla de $S$ sin coincidencia, $(w,3,0)$.

\textbf{Resultado de la join externo completo}

\begin{center}
\small
\begin{tabular}{cccc}
\toprule
\rowcolor{Fondo}
$A$ & $B$ & $C$ & $D$\\
\midrule
1 & x & 0 & 3\\
2 & y & 2 & 1\\
2 & y & 3 & 3\\
2 & y & 4 & 2\\
3 & x & 0 & 3\\
2 & z & \nulo{} & \nulo{}\\
9 & a & \nulo{} & \nulo{}\\
\nulo{} & w & 3 & 0\\
\bottomrule
\end{tabular}
\end{center}

\textbf{Cardinalidad del join: 8.} Se obtienen cinco tuplas por coincidencias y tres tuplas sin coincidencia que se conservan mediante la join externo completo.

Se agrupa el resultado por $A$ y se cuenta el número de valores de $B$ de cada grupo. Como $B$ tiene un valor en las ocho tuplas, todas contribuyen al conteo de su grupo:

\[
\gamma_{A;\ \operatorname{count}(B)\to t}(J).
\]

\begin{center}
\begin{tabular}{cc}
\toprule
\rowcolor{Fondo}
$A$ & $t=\operatorname{count}(B)$\\
\midrule
1 & 1\\
2 & 4\\
3 & 1\\
9 & 1\\
\nulo{} & 1\\
\bottomrule
\end{tabular}
\end{center}

\textbf{Cardinalidad final: 5.} Se obtiene una tupla por cada grupo distinto de $A$: 1, 2, 3, 9 y \nulo{}. El grupo $A=2$ contiene cuatro tuplas: tres provenientes de las coincidencias con $B=y$ y una correspondiente a $(2,z)$. El grupo $A=\nulo{}$ corresponde a la tupla de $S$ con $B=w$.

\clearpage

\section*{1. Evidencias de cardinalidad en RelaX}

\textbf{C08: expresión ejecutada y comprobación}

\lstinputlisting{\RutaDiego/cardinalida_c08.relax}

\begin{center}
\begin{minipage}{.48\linewidth}
\caprel{cardinalida_c08_arbol.png}{5.0}{Árbol de la expresión C08.}
\end{minipage}
\hfill
\begin{minipage}{.48\linewidth}
\caprel{cardinalidad_c08_resultado.png}{5.0}{Resultado de C08: 3 tuplas.}
\end{minipage}
\end{center}

\textbf{C10: join externo completo}

\lstinputlisting{\RutaDiego/cardinalidad_c10.relax}

\begin{center}
\begin{minipage}{.48\linewidth}
\caprel{cardinalidad_c10_arbolJoinExterno.png}{6.0}{Árbol de la join externo completo.}
\end{minipage}
\hfill
\begin{minipage}{.48\linewidth}
\caprel{cardinalidad_c10_resultadoJoinExterno.png}{6.0}{Resultado intermedio del Join: 8 tuplas.}
\end{minipage}
\end{center}

\textbf{C10: agrupación y conteo}

\begin{center}
\begin{minipage}{.48\linewidth}
\caprel{cardinalidad_c10_arbol.png}{5.0}{Árbol de la agrupación y el conteo.}
\end{minipage}
\hfill
\begin{minipage}{.48\linewidth}
\caprel{caridinalidad_c10_resultado.png}{5.0}{Resultado final de C10: 5 tuplas.}
\end{minipage}
\end{center}